\documentclass{article}
\usepackage[margin=1in]{geometry}
\usepackage{amsmath}
\usepackage{amssymb}
\usepackage{graphicx}
\usepackage{float}

\setlength{\parindent}{0pt}
\begin{document}

\begin{center}
    {\LARGE \textbf{Project}} \\[0.25cm]
    {\Large \textbf{Stacking Time-Series Momentum and Carry with Risk Budgeting}} \\[0.5cm]

    {\large Anand Nakhate, Aparna Ambarapu, Pravallika Reddy, Aman Savla, Yash Waghwala} \\[0.25cm]

    December 4, 2025
\end{center}

\vspace{0.6cm}

\begin{center}
{\Large \textbf{Abstract}}
\end{center}

This project investigates whether combining two well-established cross-asset risk premia—Time-Series
Momentum (TSMOM) and Carry—can produce a more robust, crisis-resilient futures trading strategy
when implemented with modern risk budgeting techniques. Using a diversified universe of equity index,
rates, FX, and commodity futures, we construct both back-adjusted and ratio-adjusted continuous price
series, compute trend and carry signals, and scale positions to target a consistent level of ex-ante volatility.
We then “stack” the two signals through a rolling ridge-regression framework that determines the optimal
blend between TSMOM and Carry, subject to turnover constraints and sleeve-level risk caps.\\[0\baselineskip]

Performance is evaluated across multiple regimes, including post-2015 markets, the COVID-19 shock,
and the 2022 inflation cycle. We benchmark stacked portfolios against pure TSMOM, pure Carry, equal-
weight blends, and sleeve-level risk parity. Our results show that the stacked strategy meaningfully
improves Sharpe ratio, reduces drawdowns, and exhibits lower correlation to traditional assets, indicating
that combining economically distinct signals within a risk-budgeted architecture enhances both statistical
reliability and practical implementability.

\section{Introduction}

Time-Series Momentum (TSMOM) and Carry are two of the most widely documented and economically
motivated sources of return in global futures markets. A large body of empirical research shows that
assets that have trended positively in the recent past tend to continue performing well over intermediate
horizons, while assets offering higher implied roll yield or term-structure premia generate systematically
positive excess returns. These effects have been observed consistently across equities, interest rates,
foreign exchange, and commodities, making them appealing building blocks for multi-asset quantitative
strategies.\\[0\baselineskip]

Despite their simplicity and transparency, the two signals capture fundamentally different economic
mechanisms. TSMOM is typically linked to slow-moving behavioral biases, delayed information diffusion,
and market underreaction, while Carry is associated with risk compensation for holding assets with
unbalanced supply--demand conditions, storage premia, or interest-rate differentials. Because these drivers
are distinct, TSMOM and Carry often exhibit low correlation, especially during periods of macroeconomic
stress. This creates a natural motivation to study whether combining them can produce more stable and
crisis-resilient returns than either signal alone.\\[0\baselineskip]

In practice, constructing such a combined strategy requires careful attention to several implementation
details: the formation of continuous futures price series, the treatment of roll mechanics, volatility
normalization across assets, transaction costs, and dynamic allocation between signals. Traditional
approaches often rely on fixed weighting schemes (such as 50/50 blends), which do not adapt to changing
market conditions or evolving relative performance between signals. Modern systematic portfolios,
however, increasingly make use of risk budgeting and machine-learning style regularization to determine
signal weights in a more data-driven yet stable manner.\\[0\baselineskip]

This project examines whether a risk-budgeted ``stacked'' model---which learns the optimal blend of
TSMOM and Carry via a rolling ridge regression that minimizes realized portfolio variance subject to
turnover and exposure constraints---can deliver superior out-of-sample performance. Using a diversified
universe of global futures across equities, rates, FX, and commodities, we compare the stacked strategy to
pure TSMOM, pure Carry, equal-weight combinations, and sleeve-level risk parity benchmarks. Our
analysis emphasizes robustness across recent market regimes, including the COVID-19 dislocation and
the 2022 inflation-driven macro environment.\\[0\baselineskip]

Overall, the objective of this study is twofold: first, to assess whether combining economically distinct
premia meaningfully improves portfolio performance; and second, to evaluate whether risk budgeting and
regularized stacking provide a more stable and implementable approach for multi-asset signal integration.


\section{Data and Preprocessing}

A critical component of any multi-asset systematic strategy is the construction of a clean, reliable,
and well-documented dataset. Our data pipeline integrates raw futures data sourced from Databento
into a standardized format suitable for signal generation, portfolio construction, and cross-asset
comparability. The pipeline performs loading, cleaning, metadata enrichment, sanity checks, and temporal
coverage validation before any modeling takes place.

\subsection{Universe}

We consider major exchange-traded futures across four sleeves, closely following the asset coverage in
Moskowitz et al.\ (2012) and Koijen et al.\ (2018):

\begin{itemize}
    \item \textbf{Equities}: ES (S\&P 500), NQ (Nasdaq 100), RTY (Russell 2000), NKD (Nikkei 225)
    
    \item \textbf{Rates}: ZT (2Y Treasury), ZF (5Y Treasury), ZN (10Y Treasury), 
          ZB (30Y Treasury Bond), UB (Ultra Bond)

    \item \textbf{FX (G10 currency futures)}: 6E (EUR), 6J (JPY), 6B (GBP), 
          6A (AUD), 6C (CAD), 6S (CHF), 6N (NZD)

    \item \textbf{Commodities – Energy}: CL (Crude Oil), NG (Natural Gas), HO (Heating Oil), RB (RBOB Gasoline)

    \item \textbf{Commodities – Metals}: GC (Gold), SI (Silver), HG (Copper)

    \item \textbf{Commodities – Grains}: ZC (Corn), ZW (Wheat), ZS (Soybeans), 
          ZL (Soybean Oil), ZM (Soybean Meal)
\end{itemize}

This universe mirrors typical cross-asset TSMOM and Carry implementations and provides broad coverage
across economic sectors.

\subsection{Horizon}

Where available, historical data extend back to 1990. In practice, our consolidated Databento dataset spans
January 2015 through December 2024 for most contracts, with an out-of-sample evaluation window from
2019 onward. Assets with shorter histories, such as RTY, are included but evaluated with caution.

\subsection{Fields}

For each contract, we ingest standard daily fields:

\begin{itemize}
    \item open, high, low, close, settlement prices
    \item volume and open interest
    \item contract specifications (tick size, multiplier)
    \item \texttt{front} and \texttt{back} contract prices for computing Carry
    \item metadata fields: asset class, sleeve, region, parent symbol
\end{itemize}

These fields enable continuous contract construction, roll logic, volatility estimation, and carry computation.

\subsection{Data Intake and Cleaning Pipeline}

The pipeline begins with importing the dataset and applying several cleaning steps:

\begin{enumerate}
    \item \textbf{Load raw Databento dataset}: A single consolidated CSV containing all futures symbols.
    \item \textbf{Enforce data types}: Convert date fields, numerical columns, and categorical metadata.
    \item \textbf{Sort chronologically} by contract and date.
    \item \textbf{Remove duplicate rows}: Our dataset contained 3,323 duplicates, all of which were dropped.
    \item \textbf{Metadata enrichment}: Assign each instrument its parent root, asset class (equity, rate, FX, commodity),
          sleeve, and region.
\end{enumerate}

A summary printed during execution confirms the final dataset size:

\begin{quote}
\texttt{Loaded 608,378 rows.}
\end{quote}

\subsection{Data Integrity Checks}

To ensure data quality before constructing continuous contracts or computing signals, we run a series of
sanity checks:

\begin{itemize}
    \item missing values by field
    \item negative prices (rare but possible with data vendor inconsistencies)
    \item malformed OHLC entries (e.g., high < low)
    \item checks for unrealistic or zero volumes
\end{itemize}

In our dataset, we find:

\begin{quote}
\texttt{1 row with negative prices \,(CL.FUT).}  
\texttt{0 bad OHLC rows, 0 nonpositive volume rows.}
\end{quote}

These checks allow us to flag and isolate irregularities before signal calculation.

\subsection{Universe Audit}

To validate that each parent symbol has sufficient coverage for modeling, we generate an audit table
containing:

\begin{itemize}
    \item first available date
    \item last available date
    \item number of days of data
    \item average daily contract count
    \item average daily volume
\end{itemize}

This ensures adequate historical depth for robust signal estimation and volatility scaling.

\subsection{Temporal Coverage Visualization}

Finally, we visualize the availability of each asset using a horizontal bar chart showing start and end dates.
This provides an intuitive overview of contract history and helps identify shorter-lived series such as RTY.
An example of the resulting coverage plot is shown in Figure

\begin{figure}[H]
    \centering
    \includegraphics[width=0.95\textwidth]{figures/coverage_plot.png}
    \caption{Universe Coverage: Temporal Availability across Futures Contracts.}
    \label{fig:coverage}
\end{figure}

Taken together, these preprocessing steps ensure that our dataset is clean, consistent, and ready for
constructing continuous futures series and computing TSMOM and Carry signals.

\subsection{Liquidity Profile}

In addition to temporal coverage, we assess the liquidity characteristics of each futures contract by
computing the average daily trading volume over the full sample. Liquidity is an essential consideration
for any multi-asset trading strategy, as it affects implementation costs, slippage, and the reliability of
volatility estimates.

Using the universe audit table, we sort contracts by their mean daily volume and visualize the results
with a horizontal bar chart. This allows us to quickly identify high-volume instruments such as ZN, ZF,
and ES, as well as thinner markets including HO, RB, and certain agricultural contracts.

Figure~\ref{fig:liquidity} reports the liquidity distribution across all roots in the dataset.

\begin{figure}[H]
    \centering
    \includegraphics[width=0.95\textwidth]{figures/liquidity_plot.png}
    \caption{Average Daily Volume across Futures Contracts.}
    \label{fig:liquidity}
\end{figure}

This analysis ensures that downstream modeling places appropriate emphasis on consistently traded
contracts and helps inform position scaling and turnover constraints in later stages of the strategy.


\subsection{Saving the Cleaned Dataset}

After completing all cleaning, enrichment, and integrity checks, the final dataset is exported in Parquet
format, which provides compressed columnar storage and efficient I/O for downstream analysis. The
cleaned dataset contains 608{,}378 rows across all futures roots and includes standardized fields required
for constructing continuous contracts and computing signals. The export routine writes the file
\texttt{master\_dataset\_cleaned.parquet} to the processed data directory.

\subsection{Summary of Preprocessing Pipeline}

In summary, the preprocessing workflow includes:

\begin{itemize}
    \item ingestion of raw Databento GLBX.MDP3 futures data;
    \item type enforcement, chronological sorting, and duplicate removal;
    \item metadata enrichment (asset class, sleeve, region, parent root);
    \item sanity checks for missing values, negative prices, and malformed OHLC data;
    \item universe audits for historical depth and liquidity;
    \item visualization of temporal coverage and liquidity profiles;
    \item export of a fully cleaned, analysis-ready dataset.
\end{itemize}

These steps ensure that the dataset entering the modeling stage is consistent, reliable, and suitable for
constructing continuous futures series and computing TSMOM and Carry signals.



\end{document}

